\Needspace{18\baselineskip}
\section{Concrete uniqueness of operation paths}
Choose a predicate $P$ on operation labels in a net $G$. This section first proves
a generic graph theorem; the next section shows that the single-source condition
is automatic for the typed operation labels of the linear hierarchy.

\begin{definition}[Unary operations and traversal]\label{term:unary}
Selected operations are unary when each edge whose label satisfies $P$ has a
source list of length one. An operation step $s\rightsquigarrow t$ chooses such
an edge $e$ with $s\in s(e)$ and $t\in t(e)$.
\end{definition}
\lean{Graph/OperationPaths}{17}{18}{Net.UnaryOperations}
\lean{Graph/OperationPaths}{20}{21}{Net.OperationStep}
\termdef{reverse-traversal}
Reverse traversal is the machine whose transition from $t$ to $s$ is
$s\rightsquigarrow t$ and whose successful final states are boundary inputs.
Inputs are normal: any reverse step from an input would put that input in a
target list, contradicting target exclusion.
\lean{Graph/OperationPaths}{24}{29}{Net.reverseOperations}
\begin{lemma}[Determinism derived from incidence]
If selected operations are unary, reverse operation traversal is deterministic.
\end{lemma}
\begin{proof}
Two reverse steps from the same vertex select edges both targeting that vertex.
Target uniqueness identifies the edges. Their source list has length one, so
both predecessor vertices equal its only entry and hence equal one another.
There is no additional determinism premise.
\end{proof}
\lean{Graph/OperationPaths}{31}{44}{Net.reverseOperations_deterministic}

\begin{definition}[Edge-recording trace]\label{term:operation-trace}
A reverse operation trace from $v$ to $w$ records its actual edge list. The empty
trace starts and ends at the same vertex. A nonempty trace first chooses an
operation edge targeting the current vertex and a vertex in its source, then
continues recursively from that source vertex.
\end{definition}
\lean{Graph/OperationPaths}{47}{52}{Net.OperationTrace}
A trace with edge list $l$ yields exactly $|l|$ steps in the reverse machine.
Induct on its construction: the empty trace is a zero-step execution and each
trace constructor supplies the first step of the counted execution.
\lean{Graph/OperationPaths}{54}{58}{Net.OperationTrace.steps}

\begin{theorem}[Exact uniqueness to inputs]
Under the unary condition, two reverse traces from the same vertex to inputs
have equal edge lists and equal ending input vertices.
\end{theorem}
\begin{proof}
Induct on the first trace. If it is empty, the starting vertex is an input;
the second trace cannot be nonempty, since no input belongs to a target list.
Thus it too is empty and both conclusions follow. If the first trace is
nonempty, the second cannot be empty, for the same target-exclusion reason.
Their first edges are equal by target uniqueness. Their next vertices agree by
the determinism lemma. Apply the induction hypothesis to their tails and prefix
the common edge to the resulting list equality.
\end{proof}
\lean{Graph/OperationPaths}{62}{83}{Net.OperationTrace.unique_to_inputs}
\termdef{operation-path}
A forward operation path from $s$ to $t$ with edge list $l$ means a reverse trace
from $t$ to $s$ with list $\operatorname{reverse}(l)$. Empty paths are allowed in
this endpoint-sensitive notion, unlike the nonempty edge-only path predicate of
the preceding section.
\lean{Graph/OperationPaths}{86}{87}{Net.OperationPath}
Two forward operation paths from inputs to the same vertex therefore have equal
forward edge lists and equal starting input indices. Apply exact trace uniqueness,
use injectivity of list reversal, and then injectivity of the input constructor.
\lean{Graph/OperationPaths}{89}{94}{Net.operationPath_unique_from_inputs}

\begin{theorem}[No positive cycle on an input-rooted operation path]
If a vertex $v$ is reached by a forward operation path from an input, every
reverse operation trace from $v$ back to $v$ has an empty edge list.
\end{theorem}
\begin{proof}
The input-rooted path gives a reverse execution from $v$ to a successful input.
Prefix that execution by the proposed cycle to obtain another successful
execution from $v$ to the same input. Deterministic terminal uniqueness identifies
their lengths. Thus $|\mathit{cycle}|+|\mathit{path}|=|\mathit{path}|$, so the
cycle has length zero and is empty.
\end{proof}
\lean{Graph/OperationPaths}{99}{111}{Net.no_operation_cycle_on_input_path}
This is a statement about cycles at input-reachable vertices, not a global
acyclicity restriction on nets. The self-loop example below is a valid net with
a cycle and no boundary inputs.
